% Full editable typesetting of the retained Markdown source.
% Source: INTEGRAL_EIGHT_STATE_COLLISION_DERIVATION.md
% Source SHA256: dad299495e93db162931c5ccddd86610f352ba046a9f13fb91a9fd1abea00fbe
% This conversion is not a new proof or PDF build.
% Options for packages loaded elsewhere
\PassOptionsToPackage{unicode}{hyperref}
\PassOptionsToPackage{hyphens}{url}
\documentclass[
  10pt,
]{article}
\usepackage{xcolor}
\usepackage[margin=22mm]{geometry}
\usepackage{amsmath,amssymb}
\setcounter{secnumdepth}{-\maxdimen} % remove section numbering
\usepackage{iftex}
\ifPDFTeX
  \usepackage[T1]{fontenc}
  \usepackage[utf8]{inputenc}
  \usepackage{textcomp} % provide euro and other symbols
\else % if luatex or xetex
  \usepackage{unicode-math} % this also loads fontspec
  \defaultfontfeatures{Scale=MatchLowercase}
  \defaultfontfeatures[\rmfamily]{Ligatures=TeX,Scale=1}
\fi
\usepackage{lmodern}
\ifPDFTeX\else
  % xetex/luatex font selection
\fi
% Use upquote if available, for straight quotes in verbatim environments
\IfFileExists{upquote.sty}{\usepackage{upquote}}{}
\IfFileExists{microtype.sty}{% use microtype if available
  \usepackage[]{microtype}
  \UseMicrotypeSet[protrusion]{basicmath} % disable protrusion for tt fonts
}{}
\makeatletter
\@ifundefined{KOMAClassName}{% if non-KOMA class
  \IfFileExists{parskip.sty}{%
    \usepackage{parskip}
  }{% else
    \setlength{\parindent}{0pt}
    \setlength{\parskip}{6pt plus 2pt minus 1pt}}
}{% if KOMA class
  \KOMAoptions{parskip=half}}
\makeatother
\setlength{\emergencystretch}{3em} % prevent overfull lines
\providecommand{\tightlist}{%
  \setlength{\itemsep}{0pt}\setlength{\parskip}{0pt}}
\usepackage{mathrsfs}
\usepackage{mathtools}
\usepackage{xurl}
\emergencystretch=3em
\setcounter{tocdepth}{1}
\newcommand{\Beta}{\mathrm{B}}
\DeclareUnicodeCharacter{00B2}{\ensuremath{^2}}
\DeclareUnicodeCharacter{00B1}{\ensuremath{\pm}}
\DeclareUnicodeCharacter{00D7}{\ensuremath{\times}}
\DeclareUnicodeCharacter{2192}{\ensuremath{\to}}
\DeclareUnicodeCharacter{21A6}{\ensuremath{\mapsto}}
\DeclareUnicodeCharacter{2208}{\ensuremath{\in}}
\DeclareUnicodeCharacter{2209}{\ensuremath{\notin}}
\DeclareUnicodeCharacter{2212}{\ensuremath{-}}
\DeclareUnicodeCharacter{2227}{\ensuremath{\wedge}}
\DeclareUnicodeCharacter{2228}{\ensuremath{\vee}}
\DeclareUnicodeCharacter{2260}{\ensuremath{\ne}}
\DeclareUnicodeCharacter{2264}{\ensuremath{\le}}
\DeclareUnicodeCharacter{2265}{\ensuremath{\ge}}
\DeclareUnicodeCharacter{2297}{\ensuremath{\otimes}}
\usepackage{bookmark}
\IfFileExists{xurl.sty}{\usepackage{xurl}}{} % add URL line breaks if available
\urlstyle{same}
\hypersetup{
  hidelinks,
  pdfcreator={LaTeX via pandoc}}

\author{}
\date{}

\begin{document}

\section{The integral eight-state collision and its prime-primary
classes}\label{the-integral-eight-state-collision-and-its-prime-primary-classes}

Independent derivation, 23 September 2026. The objects here are integral
coefficient algebras, cotangent complexes and ordinary de Rham
complexes. Their algebraic traces are computed explicitly. They have not
been identified with the programme's new supported-zero Weil functional.
In particular no sign conclusion for that functional is inferred from
these traces.

\subsection{1. Sources, coefficient rings and the full six-state
algebra}\label{sources-coefficient-rings-and-the-full-six-state-algebra}

The receiving programme calculations are ESH49--55 and ESH58--72 in
\texttt{EIGHT\_STATE\_HEAT\_COMPARISON.md}, read in full for this
derivation; WRF53--74 in
\texttt{WEIL\_REFLECTION\_FROBENIUS\_DEFECT\_DERIVATION.md}; and DC1--10
in \texttt{COLLISION\_DIVIDED\_CLASS\_DERIVATION.md}. The source-reading
ledger \texttt{SOURCE\_READING\_AND\_USE.md}, its collision continuation
entries, and \texttt{sources/prismatic/prisms\_index\_routes.json} were
consulted before calculation. The existing route is
PUBUNIT-3372AF0EC9F522F58C6F6CEE, Bhargav Bhatt and Peter Scholze,
\href{https://arxiv.org/abs/1905.08229v4}{Prisms and Prismatic
Cohomology, arXiv:1905.08229v4}, original \texttt{prisms.tex}, SHA256
\texttt{f155d2051758e2805efad8513690a91d7b18c19461514abdfd3d32f1676cf38b}.
The previously read source lines135--214 and484--558 supply the prism
and delta-ring definitions; the calculations below prove their
hypotheses directly. The cotangent-complex presentation is the
regular-immersion construction already sourced and applied in WRF32 and
the ledger's Illusie entry. No further literature claim or exhaustive
reading claim is made.

Put (K=\mathbb Z\_p), for a specified prime (p), and \[
B_K=K[R,T]/((R-1)^2(R+2),T^2-3R^2+3),\qquad
I_K=K[\Theta]/(\Theta^2+1).
\tag{IEC1}
\] The constant infinity algebra (I\_K) is retained as an integral
algebra. It is étale only when (p\ne2). Thus the complex infinity
chart's use of the inverse of (2\Theta) is not transferred to (p=2).
When the signed completion is interpreted through that étale chart, its
integral chart domain is (\mathbb Z{[}1/2{]}). The product
(B\_K\times I\_K) is nevertheless a specified finite free integral model
at every prime, without asserting an extension of every original affine
chart map across2.

Write (x=R-1). The two equations become \[
f=x^2(x+3),\qquad g=T^2-3x(x+2).
\tag{IEC2}
\] First divide by the monic degree-three polynomial in (x), then by the
monic degree-two polynomial in (T). This gives the unique (K)-basis \[
1,x,x^2,T,Tx,Tx^2
\tag{IEC3}
\] of (B\_K), so its rank is6. The infinity factor has basis (1,\Theta),
and the full model has rank8. These uniqueness statements hold already
over (\mathbb Z). Tensoring the resulting integral free modules with (K)
gives the displayed algebras. Their topology in this paper is the
(p)-adic topology, not a newly imposed topology in (x,T). A finite free
(K)-module is complete and separated because its coefficient sequences
are; polynomial multiplication is continuous. All maps between the
displayed finite modules given by integral matrices are consequently
continuous. Multiplication by (p) is injective.

The two equations form a regular sequence: (f) is a nonzero polynomial
in the domain (K{[}x,T{]}); modulo (f), the monic polynomial (g) is a
nonzerodivisor, since the highest (T)-coefficient of a nonzero
polynomial remains its highest coefficient after multiplication by (g).
Thus the relative cotangent complex is represented in degrees (-1,0) by
\[
L_{B_K/K}=\left[B_K\bar f\oplus B_K\bar g
\xrightarrow{J_B}B_K\,dx\oplus B_K\,dT\right],\quad
J_B=\begin{pmatrix}3x(x+2)&-6(x+1)\\0&2T\end{pmatrix}.
\tag{IEC4}
\] Here and below the matrices give images of the relation generators as
columns.

For (p\ne3), the exact root decomposition is \[
B_K\xrightarrow{\sim} C_K\times S_K,\qquad
C_K=K[\epsilon,T]/(\epsilon^2,T^2-6\epsilon),\quad
S_K=K[T_-]/(T_-^2-9),
\tag{IEC5}
\] with (R\mapsto(1+\epsilon,-2)), (T\mapsto(T,T\_-)). The idempotents
are \[
e_-=(R-1)^2/9,\qquad e_C=1-e_-.
\tag{IEC6}
\] The double-root equation and the simple-root equation differ by the
unit3. Evaluation at (R=-2) sends (e\_-) to1, while (R=1+\epsilon),
(\epsilon\^{}2=0), sends it to0. Conversely the element for a pair is
the sum of its evaluations multiplied by these two idempotents. These
substitutions prove the inverse identities, including the quadratic sign
relations. At3 the denominators in (IEC6) are unavailable. Indeed
(B\_K/p) at3 is (\mathbb F\_3{[}x,T{]}/(x\textsuperscript{3,T}2)), a
local ring, and (B\_\{\mathbb Z\_3\}) itself is local: an element with
unit scalar reduction is invertible by a convergent geometric series
after taking a power of its nilpotent reduction. A nontrivial idempotent
would give a nontrivial idempotent in its local residue ring, which is
impossible. Thus (IEC5) does not exist over (\mathbb Z\_3).

There are nonetheless surjective ring maps at every prime \[
B_K\longrightarrow C_K,quad x\mapsto\epsilon, T\mapsto T;
\qquad B_K\longrightarrow D_K=K[\eta]/\eta^2,quad x\mapsto0, T\mapsto\eta.
\tag{IEC7}
\] The first kernel is (x\textsuperscript{2B\_K=Kx}2\oplus KTx\^{}2):
the relation (x\textsuperscript{3=-3x}2) proves closure, and the four
remaining basis vectors map bijectively to the basis of (C\_K). The
second kernel has basis (x,x\textsuperscript{2,Tx,Tx}2). Both maps
preserve the reflection fixing (x) or (\epsilon) and negating the sign
coordinate. They compose through (C\_K\to D\_K).

\subsection{2. The coefficient6 defines an integral
lattice}\label{the-coefficient6-defines-an-integral-lattice}

For this section (K) may also be (\mathbb Z). The double-root algebra
has basis \[
1,\epsilon,T,\epsilon T,\qquad T^2=6\epsilon,quad T^3=6\epsilon T,quad T^4=0.
\tag{IEC8}
\] The ring map \[
K[z]/z^4\longrightarrow C_K,\qquad z\longmapsto T
\tag{IEC9}
\] is injective: the images (1,T,6\epsilon,6\epsilon T) are independent
over the torsionfree coefficient ring. Its additive cokernel is exactly
\[
(K/6K)\epsilon\oplus(K/6K)\epsilon T.
\tag{IEC10}
\] This is a statement about additive modules; (IEC9)'s image is a
unital subring, so (IEC10) is not asserted to be a ring quotient.
Inverting6 makes (IEC9) an isomorphism, with inverse
(\epsilon\mapsto z\^{}2/6). At2 and3 that inverse is unavailable and its
full failure is the two displayed primary classes.

The marked-root quotient is the square-zero extension \[
0\longrightarrow J=K\epsilon\oplus K\epsilon T
\longrightarrow C_K\xrightarrow{q}D_K\longrightarrow0,
\quad q(\epsilon)=0,quad q(T)=\eta,\quad J^2=0.
\tag{IEC11}
\] Multiplication proves (J\^{}2=0), and the bases prove exactness.
There is no (K)-algebra section of (q). Every lift of (\eta) has the
form (y=T+a\epsilon+b\epsilon T). Its square is
(6\epsilon+2a\epsilon T), which is nonzero by coefficient comparison.
Thus it cannot satisfy (y\^{}2=0). This defines a non-split square-zero
extension even though the underlying sequence of free (K)-modules
splits.

The algebraic trace is \[
\operatorname{Tr}_{C_K/K}(a+b\epsilon+cT+d\epsilon T)=4a,
\qquad\operatorname{Tr}_{D_K/K}(a+c\eta)=2a.
\tag{IEC12}
\] The nilpotent ideal (N=(\epsilon,T)) has fourth power zero.
Multiplication by each of its three basis vectors has zero diagonal in
(IEC8), while multiplication by1 has trace4. This proves the first
formula directly over (K); the second follows in the two-element basis.
In particular \[
\operatorname{Tr}_{C_K/K}=2\operatorname{Tr}_{D_K/K}\circ q.
\tag{IEC13}
\] For the reflection (\sigma(\epsilon)=\epsilon,\sigma(T)=-T), the
bilinear reflected trace is (4aa'). Its radical is exactly (N), since
(4a=0) in the torsionfree ring forces (a=0). At2 its reduction modulo2
is zero, although the integral form is not zero. These statements
concern this specified algebraic trace only.

For the whole six-state algebra write (b=P(R)+TQ(R)), with (P,Q) of
degree at most2. Polynomial multiplication or extension to the
characteristic-zero fraction field gives \[
\operatorname{Tr}_{B_K/K}(b)=4P(1)+2P(-2),
\tag{IEC14}
\] and \[
\operatorname{Tr}_{B_K/K}(\sigma(b)b')
=4P(1)P'(1)+2P(-2)P'(-2)-18Q(-2)Q'(-2).
\tag{IEC15}
\] The prime denotes the second input polynomial, not differentiation in
this formula. The characteristic-zero verification separates the root
factors using (IEC6); each expression is an integral polynomial in the
input coefficients. Since (K) embeds into its fraction field, equality
there proves equality over (K), also at3. This argument transfers an
integral identity, not the unavailable product decomposition. The
infinity factor adds (2aa'+2bb') in basis (1,\Theta) for its reflection
(\Theta\mapsto-\Theta).

\subsection{3. The exact cotangent and conormal classes of the double
root}\label{the-exact-cotangent-and-conormal-classes-of-the-double-root}

The regular-sequence proof used in (IEC4) applies to
(\epsilon\textsuperscript{2,T}2-6\epsilon). Put (f\_C=\epsilon\^{}2),
(g\_C=T\^{}2-6\epsilon). Then \[
L_{C_K/K}=\left[C_K\bar f_C\oplus C_K\bar g_C
\xrightarrow{\left(\begin{smallmatrix}2\epsilon&-6\\0&2T\end{smallmatrix}\right)}
C_Kd\epsilon\oplus C_KdT\right].
\tag{IEC16}
\] Its degree-minus-one cohomology is free of rank3 over (K), with basis
\[
\epsilon\bar f_C,\quad \epsilon T\bar f_C,\quad
3T\bar f_C+\epsilon T\bar g_C.
\tag{IEC17}
\] Indeed (T b=0) forces (b=b\_3\epsilon T), by expanding in (IEC8). The
other equation (2\epsilon a=6b) then forces the scalar coefficient of
(a) to vanish and its (T)-coefficient to equal (3b\_3). The two
remaining coefficients of (a) are unrestricted. This proves generation
and independence without dividing by2 or3.

The marked quotient's cotangent complex is \[
L_{D_K/K}=[D_K\bar h\xrightarrow{2\eta}D_Kd\eta],\qquad h=\eta^2,
\tag{IEC18}
\] with (H\^{}\{-1\}=K\eta\bar h) and
\(Omega^1_{D_K/K}=K\,d\eta\oplus(K/2K)\eta d\eta\). There is an explicit
chain map from (L\_\{C\_K/K\}\otimes\_\{C\_K\}D\_K) to (IEC18): \[
\bar f_C\mapsto0,\quad\bar g_C\mapsto\bar h,
\qquad d\epsilon\mapsto0,\quad dT\mapsto d\eta.
\tag{IEC19}
\] The two columns commute with differentials. Every vector in (IEC17)
maps to0. The transitivity connecting map to the conormal module
(J/J\^{}2=J) has the exact formula \[
K\eta\bar h\longrightarrow J,\qquad
\eta\bar h\longmapsto6\epsilon T.
\tag{IEC20}
\] To verify its coefficient directly, lift (\eta) by (T); the relation
(\eta\^{}2) lifts to (T\^{}2=6\epsilon) in (J), so its multiple by
(\eta) lifts to (6\epsilon T). This gives the indicated sign convention
for the connecting map. The ensuing conormal sequence is checked without
relying on a sign convention: \[
0\longrightarrow6K\epsilon T\longrightarrow J
\xrightarrow{d} \Omega^1_{C_K/K}\otimes_{C_K}D_K
\longrightarrow\Omega^1_{D_K/K}\longrightarrow0.
\tag{IEC21}
\] In the middle module the only nonzero relation is
(6d\epsilon-2\eta d\eta=0), together with its (\eta)-multiple. The
element ((a+b\eta)d\epsilon) vanishes exactly when (a=0) and (b\in6K):
compare both coefficients against ((u+v\eta)(6d\epsilon-2\eta d\eta)),
using torsionfreeness to deduce (u=0). This proves exactness and also
proves that (IEC20) has the stated image. Consequently the image of the
conormal map is \[
J/(6K\epsilon T)=K\epsilon\oplus(K/6K)\epsilon T.
\tag{IEC22}
\] Its second factor is a genuine2- or3-primary cotangent class at the
corresponding prime.

\subsection{4. Ordinary de Rham cohomology, including degree
two}\label{ordinary-de-rham-cohomology-including-degree-two}

All complexes in this section are ordinary relative Kähler de Rham
complexes, with (Omega\textsuperscript{i=\bigwedge}i\Omega\^{}1). They
are finite (K)-modules and already (p)-adically complete. They are not
silently substituted for the divided-power complex used in CPD/DC.

Put (a=d\epsilon), (b=dT), (z=a\wedge b). Expanding the two
differentiated relations and multiplying by the basis (IEC8) gives the
complete (K)-module presentation \[
\begin{gathered}
6a-2Tb=0,\quad 2\epsilon a=0,\quad2\epsilon T a=0,
\quad2\epsilon T b=0,\quad6(Ta-2\epsilon b)=0.
\end{gathered}
\tag{IEC23}
\] These generate all relation multiples: multiples of (2\epsilon a) by
(1,\epsilon,T,\epsilon T) give its two displayed nonzero multiples;
multiples of (6a-2Tb) give itself, (6\epsilon a-2\epsilon Tb),
(6Ta-12\epsilon b), and (6\epsilon Ta). The first relation family
reduces the latter three exactly to those in (IEC23). Thus no relation
was dropped.

In particular \[
\Omega^1_{C_K/K}\simeq K^3\oplus(K/6K)\oplus(K/2K)^4
\tag{IEC24}
\] as (K)-modules: the ((a,Tb)) relation has coefficient gcd2, the
((Ta,\epsilon b)) relation has coefficient gcd6, and the three remaining
relations have coefficient2. Each two-generator one-relation block has
one free generator and the indicated torsion; (b) supplies the third
free generator. This proof also works over (\mathbb Z).

The exterior relation module gives \[
\Omega^2_{C_K/K}=C_K/(6,2\epsilon,2T)\,z,
\qquad\Omega^i=0\quad(i>2).
\tag{IEC25}
\] Wedging (2\epsilon a) with (b) gives (2\epsilon z); wedging (2Tb-6a)
with (a) and with (b) gives (2Tz) and (6z). Conversely these generate
all wedge relations, proving the displayed quotient.

The degree-zero derivative has kernel exactly (K). One may check this
after embedding into the fraction field, where (C) is
(K{[}1/p{]}{[}T{]}/T\^{}4) and the derivatives of
(T,T\textsuperscript{2,T}3) are independent in characteristic0.
Alternatively insert
(d(A+B\epsilon+CT+D\epsilon T)=Ba+Cb+D(Ta+\epsilon b)) into (IEC23); its
free components force (B=C=D=0).

To compute the rest, first quotient (Omega\^{}1) by (dC\_K). Setting
(a=b=0) and (Ta=-\epsilon b) in the additive quotient yields \[
\Omega^1/dC_K=(K/18K)v\oplus(K/2K)u\oplus(K/2K)q_0
\oplus(K/2K)r_0\oplus(K/2K)s_0,
\tag{IEC26}
\] where (v=\epsilon dT), (u=T dT), (q\_0=\epsilon d\epsilon),
(r\_0=\epsilon T d\epsilon), (s\_0=\epsilon T dT). The coefficient18 is
(6\cdot3), coming from (6(Ta-2\epsilon b)) after (Ta=-\epsilon b). The
induced differential is exactly \[
v\mapsto z,\quad u,q_0\mapsto0,\quad
r_0\mapsto-\epsilon z,\quad s_0\mapsto Tz.
\tag{IEC27}
\] The coefficients of (z,\epsilon z,Tz,\epsilon Tz) in (IEC25) have
orders6,2,2,2 over (\mathbb Z), respectively. Equations (IEC26)--(IEC27)
therefore give \[
\begin{array}{c|c|c|c}
K&H^0&H^1&H^2\\\hline
\mathbb Z&\mathbb Z&(\mathbb Z/2)[\epsilon d\epsilon]\oplus(\mathbb Z/2)[T dT]\oplus(\mathbb Z/3)[6\epsilon dT]&(\mathbb Z/2)[\epsilon Tz]\\
\mathbb Z_2&K&\mathbb F_2[\epsilon d\epsilon]\oplus\mathbb F_2[T dT]&\mathbb F_2[\epsilon Tz]\\
\mathbb Z_3&K&\mathbb F_3[3\epsilon dT]&0\\
\mathbb Z_p,\ p>3&K&0&0.
\end{array}
\tag{IEC28}
\] Each displayed generator has the indicated exact order: (IEC26) is a
direct cyclic decomposition, and (IEC27) determines its kernel
coefficient by coefficient. For example at3 the (v)-summand has order9
and maps onto an order3 copy, leaving (3v) of order3. At2 that same
summand maps injectively onto the scalar (z)-summand; the two additional
mixed generators map injectively onto (\epsilon z,Tz), leaving exactly
(u,q\_0). No generator maps to (\epsilon Tz), proving the degree-two
class.

Under (q:C\_K\to D\_K), \[
[T dT]\longmapsto[\eta d\eta],\qquad
[\epsilon d\epsilon],\ [3\epsilon dT],\ [\epsilon Tz]\longmapsto0.
\tag{IEC29}
\] This follows from (d\epsilon\mapsto0,T\mapsto\eta), and the target de
Rham group is (H\^{}1(D\_K)=(K/2K){[}\eta d\eta{]}). Thus the map is
onto the earlier WRF/DC collision class at2; its other classes are
retained in the stated kernel. At3 the nonzero ordinary class is killed
by this quotient.

The second order-two class also has an exact dual-number source: the
injective ring map (j\_\epsilon:D\_K\to C\_K), (\eta\mapsto\epsilon),
sends ({[}\eta d\eta{]}) to ({[}\epsilon d\epsilon{]}). Injectivity
follows from the independent basis vectors (1,\epsilon), and the
relation is preserved since (\epsilon\^{}2=0). With the monoidal lifts
this is a delta and prism map. It is not a reflection-equivariant map
for the previously specified sign reflection on (D\_K): its generator
discrepancy is (2\epsilon), because (\epsilon) is reflection-even. On
cohomology at2 the scalar discrepancy vanishes, so both degree-one
classes are reflection-invariant. Their product is exactly \[
[\epsilon d\epsilon]\,[TdT]=[\epsilon T d\epsilon\wedge dT]\ne0.
\tag{IEC60}
\] The degree-two nonvanishing was proved in (IEC28), and each generator
squares to0 because (d\epsilon\wedge d\epsilon=dT\wedge dT=0). Thus the
full cohomology multiplication at2 is specified by (H\^{}0=K), the two
order-two degree-one generators, their one order-two product, and zero
higher products. At3 the class ({[}3\epsilon dT{]}) is reflection-odd.
These are actions of the explicit algebraic reflection, not an
identification with an analytic holonomy operator.

\subsection{5. Frobenius lifts on the double-root
lattice}\label{frobenius-lifts-on-the-double-root-lattice}

For every prime the monoidal formulas \[
\Phi^{\mathrm m}(\epsilon)=0,\quad \Phi^{\mathrm m}(T)=T^p,
\quad\Phi^{\mathrm m}|K=\mathrm{id}
\tag{IEC30}
\] define an endomorphism of (C\_K). The relation image is
(T\textsuperscript{\{2p\}=6}p\epsilon\^{}p=0), while the image of
(6\epsilon) is0. The map reduces to Frobenius modulo (p), including on
coefficients since (a\^{}p\equiv a\pmod p) in (K). Thus
\(\delta^{\mathrm m}(b)=(\Phi^{\mathrm m}(b)-b^p)/p\) is a well-defined
delta operation. Its addition and multiplication identities follow by
expanding the numerator; torsionfreeness gives uniqueness. In particular
(\delta\^{}\{\mathrm m\}(\epsilon)=\delta\^{}\{\mathrm m\}(T)=0).

There is also the augmentation lift \[
\Phi^0(\epsilon)=\Phi^0(T)=0,\qquad
\delta^0(T)=-T^p/p,quad\delta^0(\epsilon)=0.
\tag{IEC31}
\] Indeed (T\^{}2=6\epsilon) is divisible by2, (T\^{}3=6\epsilon T) is
divisible by3, and (T\^{}p=0) for (p\ge5). Hence the augmentation
reduces to Frobenius at every prime. Explicitly
(\delta\^{}0(T)=-3\epsilon) at2 and (=-2\epsilon T) at3. It commutes
with the reflection, so its delta operation does too.

For completeness all coefficient-fixed Frobenius lifts on (C\_K) are \[
\Phi_v(T)=p v,\qquad \Phi_v(\epsilon)=p^2v^2/6,
\qquad v=b\epsilon+cT+d\epsilon T\in N.
\tag{IEC32}
\] Here the notation for the second formula means the integral
expression \[
p^2v^2/6=p^2c^2\epsilon+(p^2bc/3)\epsilon T.
\tag{IEC33}
\] At3, (p\^{}2/3=3); at other primes3 is a unit, so every coefficient
is defined. A nilpotent must have zero scalar coefficient in (IEC8). A
ring endomorphism sends the nilpotent (T) to a nilpotent; its Frobenius
congruence requires its three coefficients to be divisible by (p), since
(T\^{}p\in pC\_K). Hence it is (pv). The relation
(6\Phi(\epsilon)=\Phi(T)\^{}2), in a torsionfree ring, forces (IEC33).
Conversely (IEC33) squares to0 and its product with6 equals ((pv)\^{}2).
Its coefficients are divisible by (p), proving the required Frobenius
congruence for (\epsilon). These facts prove existence and exhaustion.

The lift (IEC32) is reflection-equivariant exactly when (b=0). Necessity
follows by comparing (\sigma\Phi(T)) with (-\Phi(T)): their difference
is (2pb\epsilon), zero only when (b=0). For (b=0),
(\Phi(\epsilon)=p\textsuperscript{2c}2\epsilon) is invariant and
(\Phi(T)) is anti-invariant, so the equality holds on both generators
and hence on the full ring. All such lifts descend along (q), sending
(\eta\mapsto pc\eta), with (\delta(\eta)=c\eta). The previous WRF
collision lift is the case (c=0). Thus not every lift is asserted to
descend to the same delta structure on (D\_K).

For the monoidal lift at2 and a general element
(A+B\epsilon+CT+D\epsilon T), the full equivariance defects are \[
(\sigma\Phi^{\mathrm m}-\Phi^{\mathrm m}\sigma)(b)=12C\epsilon,
\qquad
(\sigma\delta^{\mathrm m}-\delta^{\mathrm m}\sigma)(b)=6C\epsilon.
\tag{IEC34}
\] The first follows by applying (IEC30), which kills (\epsilon) and
(\epsilon T) and sends (T) to (6\epsilon). For the second, the ordinary
square commutes with the ring automorphism (\sigma), so the two delta
numerators differ by the first defect and can be divided by2. At odd
primes the monoidal lift commutes with the reflection because (T\^{}p)
is anti-invariant; its delta operation therefore commutes as well.

The universal ideals killing the two defects at2 are respectively \[
(12\epsilon),\qquad(6\epsilon).
\tag{IEC35}
\] Every defect lies in the indicated ideal and its value on (T) is the
generator, proving necessity and sufficiency for every quotient ring
with descended operations. Both ideals are already delta-stable: every
element in either is square zero and its monoidal Frobenius is0, so its
delta is0. They are also reflection-stable. Their full additive
quotients are \[
\begin{aligned}
Q_\Phi&=K\{1,T\}\oplus(K/4K)\{\epsilon,\epsilon T\},\quad T^2=6\epsilon,\\
Q_\delta&=K\{1,T\}\oplus(K/2K)\{\epsilon,\epsilon T\},\quad T^2=0.
\end{aligned}
\tag{IEC36}
\] Here3 is a unit in (K=\mathbb Z\_2); no class in either torsion
summand is discarded. The first quotient makes Frobenius equivariant but
leaves the delta defect (6\epsilon\ne0). The second kills that delta
defect. Their local maximal ideal is generated by2, (\epsilon,T). A
common nonzero annihilator of this maximal ideal is (2\epsilon T) in
(Q\_\Phi), and (\epsilon T) in (Q\_\delta). Thus every nonunit is a zero
divisor. A Cartier ideal in a local ring must be generated by a
nonzerodivisor, so it would have to be the unit ideal. Such an ideal
cannot be a prism ideal on these nonzero rings, since completion for
((p,I)=(1)) is zero. Hence neither quotient admits a prism ideal.

In contrast, for every lift (IEC32), the pair ((C\_K,(p))) is a bounded
crystalline prism: (C\_K) is (p)-torsionfree and (p)-complete by (IEC8),
((p)) is Cartier, and (p\in(p,\Phi(p))); the quotient has (p)-torsion
killed by (p). The same verification applies to (D\_K) for each
indicated descended lift. The maps are prism maps exactly when their
specified Frobenius operations commute; in a torsionfree target this
also gives delta compatibility by dividing the equality by (p).

\subsection{6. Exact comparison to the WRF collision and the
divided-power
class}\label{exact-comparison-to-the-wrf-collision-and-the-divided-power-class}

Take the common WRF coefficient model \[
E_K=K[[u]][r]/(r^2-u),\qquad
u\longmapsto6\epsilon,\quad r\longmapsto T.
\tag{IEC37}
\] Evaluation of a formal series is well-defined because
(\epsilon\^{}2=0): only its first two (u)-coefficients contribute. The
kernel is exactly ((u\^{}2)), as the independent images
(1,T,6\epsilon,6\epsilon T) show. Thus its image and additive cokernel
are precisely (IEC9)--(IEC10). With the WRF monoidal lift
(\Phi(u)=u\textsuperscript{p,\Phi(r)=r}p), and (IEC30) on the target,
this is a Frobenius and delta map: both images of (u) under the
Frobenius square are0, and both images of (r) are (T\^{}p). It also
commutes with the reflections. Composing with (q) gives
(u\mapsto0,r\mapsto\eta), the original WRF collision quotient with its
exact sign generator. No rescaling by (2\sqrt2) and no map from complex
scalars to (K) is used.

For the augmentation lift on (C\_K), the same underlying ring map is
generally not a delta map from this specified WRF lift. Its Frobenius
discrepancy on (r) is (-T\^{}p), so its delta discrepancy is
(-T\^{}p/p): at2 it is (-3\epsilon), at3 it is (-2\epsilon T), and at
(p\textgreater3) it is0. This computes the defect exactly, instead of
assigning compatibility from the shared underlying ring map.

The PD class computed in DC5 has source (\widehat P), the completed
divided-power envelope of (z\^{}2), and generator
(\theta\emph{p={[}z\^{}\{p-1\}dz{]}). Write its complete restricted
basis as (E\_n=z\^{}\{2n\}/n!), (O\_n=z\^{}\{2n+1\}/n!), with
coefficients tending to zero (p)-adically. There is a continuous ring
map \[
\rho:\widehat P\longrightarrow C_K,\quad
E_0\mapsto1,\ O_0\mapsto T,\ E_1\mapsto6\epsilon,
\ O_1\mapsto6\epsilon T,\ E_n,O_n\mapsto0\quad(n\ge2).
\tag{IEC38}
\] The polynomial evaluation (z\mapsto T) into (C\_K{[}1/p{]}) gives
exactly these images, because (T\^{}4=0). They are all integral, so
evaluation restricts to a ring map on the uncompleted PD envelope. Only
four coefficients contribute, proving continuity and a unique extension
to the completion. Equivalently direct substitution in
(E\_iE\_j=\binom{i+j}{i}E}\{i+j\}), (E\_iO\_j=\binom{i+j}{i}O\_\{i+j\}),
and (O\_iO\_j=(i+j+1)!/(i!j!)E\_\{i+j+1\}) proves the same
multiplication statement without a rational target. Its image is the
subring (IEC9), with the entire additive cokernel (IEC10).

It extends to the exact cochain map from the two-term PD de Rham complex
to the ordinary de Rham complex of (C\_K), by (g,dz\mapsto\rho(g)dT). To
check the PD differential identities in the torsion-bearing target, the
only nonautomatic cases are \[
d(6\epsilon)=2T\,dT,\quad
d(6\epsilon T)=18\epsilon\,dT,\quad
0=12\epsilon T\,dT.
\tag{IEC39}
\] The first is the derivative of (T\^{}2=6\epsilon), the second follows
from (6T,d\epsilon=12\epsilon,dT) in (IEC23), and the third follows from
(2\epsilon T,dT=0). They correspond respectively to (dE\_1=2O\_0dz),
(dO\_1=3E\_1dz), and (dE\_2=2O\_1dz). All higher terms map to0 on both
sides, and (dO\_0=dz) is immediate. Finally
(d(\rho(g)dT)=d\rho(g)\wedge dT=\rho(dg/dz)dT\wedge dT=0), so the map
also has the required zero degree-two source component.

Consequently the complete coefficient-preserving comparison is \[
\widehat P\xrightarrow{\rho}C_K\xrightarrow{q}D_K,
\quad
\theta_2\longmapsto[TdT]\longmapsto[\eta d\eta]\ne0,
\quad
\theta_3\longmapsto[6\epsilon dT]\ne0\longmapsto0.
\tag{IEC40}
\] The first row is at2 and the second at3. At3, (6\epsilon dT) is twice
the generator in (IEC28), so it has exact order3. For (p\textgreater3),
(T\^{}\{p-1\}=0), so (\theta\_p) maps to0. The additional class
({[}\epsilon d\epsilon{]}) at2 is not in the image of the PD (H\^{}1):
DC10 gives basis classes ({[}O\_n dz{]}); only (n=0) survives, since
(O\_1dT) maps to (6\epsilon T dT=0) and all later terms map to0. The
degree-two class has no degree-two source. Thus both the surviving PD
classes and the extra integral-lattice classes are explicitly retained.

\subsection{7. The unsplit algebra at3 and its actual
maps}\label{the-unsplit-algebra-at3-and-its-actual-maps}

The complete independent calculation at3 is included below as Appendix
A, with proof labels P3.1--P3.14. It proves every module decomposition
used in this section, including the elimination of all cotangent
relations. With (n\_1=x(x+3)), (n\_2=T(x+3)), (n\_3=Tx(x+3)), its
nilradical and powers are \[
N_B=Kn_1\oplus Kn_2\oplus Kn_3,\quad
N_B^2=18Kn_1\oplus3Kn_3,\quad
N_B^3=54Kn_3,\quad N_B^4=0.
\tag{IEC41}
\] Its reduced quotient is the order
(\{(a,b,c)\in K\^{}3:a\equiv b\equiv c\pmod3\}), with conductor
(3K\^{}3) in (K\^{}3). More strongly, the coordinate-preserving
injection \[
B_K\longrightarrow C_K\times K\times K,
\quad x\mapsto(\epsilon,-3,-3),\quad T\mapsto(T,3,-3)
\tag{IEC42}
\] has additive cokernel (K/9K\oplus K/27K). Appendix A, P3.6 gives both
congruences defining its image, so this does not suppress the nonsplit
integral gluing. The trace (IEC14) has weights4,1,1 on the three factors
in (IEC42); summing the last two weights recovers the simple-root
factor's weight2.

The complete ordinary de Rham groups at3 are \[
\begin{aligned}
H^0(B_K/K)&=K\cdot1\oplus K\cdot3x^2\oplus K\cdot3x^2T,\\
H^1(B_K/K)&=\mathbb F_3[x^2dx]\oplus\mathbb F_3[3x\,dT],\\
H^j(B_K/K)&=0\quad(j\ge2).
\end{aligned}
\tag{IEC43}
\] Under the first quotient in (IEC7), the degree-one map fits into the
split exact sequence of additive groups \[
0\longrightarrow\mathbb F_3[x^2dx]
\longrightarrow H^1(B_K/K)
\longrightarrow H^1(C_K/K)\longrightarrow0,
\quad[3x\,dT]\longmapsto[3\epsilon\,dT].
\tag{IEC44}
\] Substitution proves these images, and the independent generators in
(IEC28), (IEC43) prove exactness. Both classes are killed by the
subsequent dual-number quotient. There is no assertion that the
splitting of this sequence supplies an algebra section: (IEC11) and
Appendix A, P3.5 prove the failure of those algebra sections.

The map of cotangent complexes (B\_K\to C\_K) sends \[
\bar f\mapsto(\epsilon+3)\bar f_C,
\quad\bar g\mapsto-3\bar f_C+\bar g_C,
\quad dx\mapsto d\epsilon,quad dT\mapsto dT.
\tag{IEC45}
\] Indeed (f(\epsilon)=\epsilon\^{}2(\epsilon+3)) and
(g(\epsilon,T)=g\_C-3f\_C); differentiating gives the degree-zero
square. Appendix A gives the basis of (H\^{}\{-1\}(L\_\{B\_K/K\})) as
((n\_1,0),(n\_3,0),(n\_2,n\_3)). In the basis
(v\_1=\epsilon\bar f\_C,v\_2=\epsilon T\bar f\_C,v\_3=3T\bar f\_C+\epsilon T\bar g\_C)
from (IEC17), the images are \[
9v_1,\qquad9v_2,\qquad3v_3-3v_2.
\tag{IEC46}
\] For example (n\_2\mapsto3T+\epsilon T), (n\_3\mapsto3\epsilon T);
substitution in (IEC45) gives
(9T\bar f\_C-3\epsilon T\bar f\_C+3\epsilon T\bar g\_C), the third
expression. The map is injective, with cokernel \[
(K/9K)^2\oplus K/3K.
\tag{IEC47}
\] Replace (v\_3) by (v\_3-v\_2) to diagonalize (IEC46); this is an
integral basis change. Thus the cotangent lattice itself detects both
order9 and order3 classes. Its further map to
(H\^{}\{-1\}(L\_\{D\_K/K\})) is zero by (IEC19).

There is a direct continuous lift of the PD comparison to the whole
(B\_\{\mathbb Z\_3\}): \[
\rho_B(E_n)=\frac{3^n}{n!}[x(x+2)]^n,
\qquad \rho_B(O_n)=T\rho_B(E_n).
\tag{IEC48}
\] The fraction (3\^{}n/n!) is integral because
(v\_3(n!)=\sum\_\{j\ge1\}\lfloor n/3\^{}j\rfloor\le n/2). Evaluation
(z\mapsto T) in the fraction-field algebra gives these formulas, since
(T\^{}2=3x(x+2)); it proves all multiplication identities, and the
images lie in the integral ring. Their valuations are bounded below by0.
Therefore restricted coefficient sums converge in the finite free
complete target, proving continuity and extension to (\widehat P). The
target differential identities follow without division in its torsion
module: writing (U=x(x+2)), \[
d\rho_B(E_n)=\frac{2\cdot3^n}{(n-1)!}U^{n-1}(x+1)dx
=2\rho_B(O_{n-1})dT,
\tag{IEC49}
\] where the second equality uses (2TdT=6(x+1)dx) multiplied by the
integral coefficient (3\^{}\{n-1\}/(n-1)!). Then Leibniz and (T\^{}2=3U)
give (d\rho\_B(O\_n)=(2n+1)\rho\_B(E\_n)dT). These identities include
(n=0) in the latter formula, while the former starts at (n=1). They
prove the cochain map to ordinary de Rham and its continuous extension.
Composition with (B\_K\to C\_K) is (IEC38): (U\mapsto2\epsilon), so its
square vanishes.

The exact retained cohomology path at3 is therefore \[
\theta_3\longmapsto[T^2dT]=[6x\,dT]
\longmapsto[6\epsilon\,dT]\longmapsto0.
\tag{IEC50}
\] In the first equality (T\textsuperscript{2dT=3x}2dT+6x,dT), and
(3x\^{}2dT=0) by Appendix A's complete cotangent presentation. The first
two nonzero images have exact order3 by (IEC43) and (IEC28). Thus this
class survives in the full unsplit six-state receiver before its final
quotient.

All coefficient-fixed Frobenius lifts on (B\_\{\mathbb Z\_3\}) are, as
proved exhaustively in Appendix A, P3.7, \[
(x,T)\mapsto(-3,3),\quad(x,T)\mapsto(-3,-3),\quad
\Phi_v(x)=\frac32v^2,\quad\Phi_v(T)=3v\quad(v\in N_B).
\tag{IEC51}
\] The augmentation (v=0) is reflection-equivariant and commutes with
the maps (B\_K\to C\_K\to D\_K) when those rings have augmentation
lifts. Its delta values are (\delta(x)=x\^{}2), (\delta(T)=-Tx(x+2)); on
(C\_K) they become (0,-2\epsilon T), exactly (IEC31). These are maps of
bounded crystalline prisms with ideal ((3)), by the finite-free
completeness argument after (IEC36).

The canonical monoidal PD lift (z\mapsto z\^{}3) is not compatible with
(IEC48) for any lift in (IEC51). It would require (\Phi(T)=T\^{}3). For
the two constant lifts, evaluate at ((x,T)=(0,0)): their values are
(\pm3), whereas (T\^{}3) evaluates to0. For (\Phi\_v), equality would
force (v=Tx(x+2)). Evaluation at ((-3,3)) gives9, while every nilpotent
(v\in N\_B) has evaluation0. Hence equality is impossible. The exact
discrepancy on the specified generator is \[
\Phi_v\rho_B(z)-\rho_B\Phi_{\widehat P}(z)
=3v-T^3,
\quad
\delta_v\rho_B(z)-\rho_B\delta_{\widehat P}(z)
=v-Tx(x+2).
\tag{IEC52}
\] This does not negate (IEC48)--(IEC50), which are genuine ring and
cochain maps. It identifies the extra Frobenius data those maps do not
preserve. Every lift in (IEC51) nevertheless induces the zero map on the
ordinary degree-one groups (IEC43): Appendix A, P3.14 provides the
actual primitives. Thus survival of the class and its ordinary Frobenius
image are different computed questions.

The same vanishing holds for all the lifts (IEC32) on positive-degree
ordinary cohomology of (C\_K). At2 each generator of (H\^{}1) has form
(y,dy), where (y=\epsilon) or (T), and its image is (4v,dv=2d(v\^{}2))
because the image of (y) is divisible by2. The degree-two image is
divisible by2 in (Omega\^{}2), which is killed by2, so it is zero. At3
the image of (3\epsilon dT) is
(\frac{27}{2}v\textsuperscript{2dv=\frac92d(v}3)), again exact.
Higher-prime groups vanish by (IEC28). These are actual primitives, not
deductions from vanishing after rationalization.

The original-coordinate monoidal candidate
(R\mapsto R\^{}3,T\mapsto T\^{}3) does not itself descend to
(B\_\{\mathbb Z\_3\}). Its full relation defects, retaining (x=R-1), are
\[
\Phi(x)=3x,\qquad
\Phi(f)=-54x^2,\qquad \Phi(g)=54x^2-18x
\quad\text{in }B_K.
\tag{IEC56}
\] Expand ((1+x)\^{}3-1) and use (x\textsuperscript{3=-3x}2) to get its
first value; substitution in (f,g) gives the next two. The ideal
generated by the relation defects is exactly ((18x)=(9x)). It is stable
under the candidate since (\Phi(18x)=54x), so the candidate does give an
endomorphism of that quotient, with all its root-coordinate torsion
retained.

Its universal delta quotient is stronger. Work in the (3)-adic
completion of the polynomial ring (K{[}x,T{]}), with the actual delta
structure coming from (\Phi(x)=(1+x)\^{}3-1), (\Phi(T)=T\^{}3). This
completed polynomial ring is torsionfree; its Frobenius lift is
continuous and reduces to Frobenius. The smallest closed delta-stable
ideal containing (f,g) is \[
(f,g,3x).
\tag{IEC57}
\] Indeed the images of (\delta(f),\delta(g)) modulo ((f,g)) are
(-18x\^{}2) and (18x\^{}2-6x=6x(3x-1)), by dividing (IEC56) by3. The
element (3x-1) is a unit in the completed polynomial ring, with inverse
(-\sum\_\{n\ge0\}(3x)\^{}n). Thus every such ideal must contain (3x).
Conversely these two delta images lie in ((3x)), while \[
\delta(x)=x+x^2,\qquad
\delta(3x)=3x+27x^2\quad\bmod(f,g).
\tag{IEC58}
\] The first follows directly from (Phi(x)); the second follows by
expanding ((\Phi(3x)-(3x)\^{}3)/3) and reducing
(x\textsuperscript{3=-3x}2). Both lie in the required ideal when
appropriate. The delta product and addition identities then prove
stability of the ideal generated by (f,g,3x): generator multiples have
delta in the ideal, and the cross terms in delta of a sum are products
of ideal elements. The quotient is finite over (K), hence complete, so
the ideal is closed. Continuity preserves its closure. This proves both
directions of the universal assertion.

The retained obstruction ring and its exact collision map are \[
Q_3=B_K/(3x)
=K\{1,T\}\oplus\mathbb F_3\{x,x^2,Tx,Tx^2\},
\quad x^3=T^2=0,
\quad Q_3\longrightarrow D_K, x\mapsto0, T\mapsto\eta.
\tag{IEC59}
\] Its induced Frobenius kills (x,T), but its induced delta is
(\delta(x)=x+x\^{}2), (\delta(T)=0); a delta operation on a ring
with3-torsion cannot be reconstructed by choosing a quotient of
((\Phi-b\^{}3)/3). The four displayed torsion basis vectors form the
full kernel of the map to (D\_K), and that map is a delta map by these
generator formulas and the quotient construction. The ring (Q\_3) has no
prism ideal: it is local with maximal ideal ((3,x,T)), and its nonzero
socle element (x\^{}2T) is annihilated by all three generators. The same
Cartier-ideal argument as in (IEC36) applies. The torsionfree
augmentation prism on (B\_K) and this universal monoidal delta quotient
are therefore two different exact constructions, with their stated maps
and retained classes.

At2, a map (\widehat P\to B\_\{\mathbb Z\_2\}) sending (z) to the whole
sign coordinate (T) cannot exist. Its generator (E\_2=z\^{}4/2) would
require an element with twice its image equal to (T\^{}4). Projection to
(S\_K) in (IEC5) gives (T\_-\^{}4=81), which is not divisible by2. This
is an exact obstruction. The stable alternative is the unital map to
(C\_K\times S\_K) given by ((\rho,\operatorname{aug})), with
(operatorname\{aug\}(z)=0); under (IEC5), its (z)-image is (e\_C T). It
carries ( heta\_2) to the class supported on (C\_K) and retains the
separate (S\_K) factor. An augmentation into (I\_K) likewise extends it
to the specified rank-eight integral model. This is an actual map with
its exact projected sign coordinate, not a claim that (e\_CT=T).

\subsection{8. The simple-root and infinity factors
at2}\label{the-simple-root-and-infinity-factors-at2}

For (A=K{[}t{]}/(t\^{}2-a)), with (a\in K\^{}\times) at2, (t) is a unit
and \[
L_{A/K}=[A\bar h\xrightarrow{2t}A\,dt],\qquad
H^{-1}=0,\quad\Omega^1=A/2A\,dt.
\tag{IEC53}
\] The nonzerodivisor equation gives the displayed complex, and (2t) is
a nonzerodivisor since (A) is free and (t) a unit. Differentiating
(b+ct) gives (c,dt), zero exactly when (c\in2K). The exterior square is
zero because (Omega\^{}1) has one generator. Hence \[
H^0(A/K)=K\oplus2Kt,\qquad
H^1(A/K)=\mathbb F_2[t\,dt],\qquad H^{\ge2}=0.
\tag{IEC54}
\] This applies to (S\_\{\mathbb Z\_2\}), with (a=9), and to
(I\_\{\mathbb Z\_2\}), with (a=-1). At odd primes (I\_K) is étale:
(2\Theta) is a unit, so its cotangent complex is acyclic,
(Omega\^{}1=0), and its de Rham complex is (I\_K) in degree0. For (S\_K)
at primes other than2 and3, (2T\_-) is a unit and the same conclusion
applies; the two signs split using the idempotents ((1\pm T\_-/3)/2),
whose denominators are units precisely there.

Finite products split Kähler differentials and their exterior complexes
componentwise. Indeed a derivation kills an idempotent (e): (de=2e,de),
and (1-2e) is a unit because its square is1. Every derivation and every
wedge product therefore separates across the factors. Equations (IEC5),
(IEC28), (IEC54) give \[
\begin{aligned}
H^1(B_{\mathbb Z_2}/K)&=\mathbb F_2[\epsilon d\epsilon]
\oplus\mathbb F_2[TdT]_C\oplus\mathbb F_2[T_-dT_-]_S,\\
H^2(B_{\mathbb Z_2}/K)&=\mathbb F_2[\epsilon T d\epsilon\wedge dT]_C,\\
H^1((B_{\mathbb Z_2}\times I_{\mathbb Z_2})/K)
&=H^1(B_{\mathbb Z_2}/K)\oplus\mathbb F_2[\Theta d\Theta],
\end{aligned}
\tag{IEC55}
\] and adjoining the infinity factor leaves the degree-two group
unchanged. The subscript (C) specifies support on the double-root
factor, not a scalar field. The degree-zero groups are the corresponding
products of (K), (K\oplus2KT\_-), and (K\oplus2K\Theta). At3 the
infinity factor adds no positive-degree class, so the full rank-eight
model has precisely the two positive-degree classes (IEC43).

The simple-root ring (S\_\{\mathbb Z\_2\}) has Frobenius lifts
(T\_-\mapsto3) and (T\_-\mapsto-3). They respect its relation and reduce
to Frobenius because (T\_-\^{}2\equiv1\pmod2). There is no
reflection-equivariant Frobenius lift on this ring. To prove exhaustion,
embed it in (\mathbb Q\_2\times\mathbb Q\_2) by evaluations at
(T\_-=\pm3). An element squaring to9 has coordinates chosen
independently from (3,-3); the four resulting elements are
(3,-3,T\_-,-T\_-), all integral. The last two do not reduce to
Frobenius: (T\_-\ne1) in (S\_K/2=\mathbb F\_2{[}T\_-{]}/(T\_--1)\^{}2).
The first two fail reflection-equivariance because the reflection fixes
constants and negates the source generator. Thus these two constant
lifts are all the lifts. Together with any chosen lift on (C\_K), they
give a Frobenius lift and a crystalline prism on (B\_\{\mathbb Z\_2\})
through (IEC5); none of these whole-ring lifts is
reflection-equivariant. In fact any endomorphism lifting Frobenius fixes
its two idempotents, since it fixes their reductions and idempotents
have unique lifts across the (2)-adically complete ideal ((2)). For the
latter uniqueness, if commuting idempotents (e,f) have the same
reduction, the idempotents (e(1-f)), (f(1-e)) belong to (2B), and an
idempotent in (2B) is divisible by every (2\^{}n), hence zero. The
obstruction from (S\_K) consequently rules out every
reflection-equivariant lift on (B\_\{\mathbb Z\_2\}).

There is no Frobenius lift at all on (I\_\{\mathbb Z\_2\}). If an image
(y=a+b\Theta) has (y\^{}2=-1), its coefficients satisfy (2ab=0) and
(a\textsuperscript{2-b}2=-1). Thus either (b=0), which would give
(a\^{}2=-1) in (\mathbb Z\_2), or (a=0,b=\pm1). The first case is
impossible modulo4, since a square is0 or1; the latter gives
(y=\pm\Theta). Modulo2 these both equal (\Theta), whereas absolute
Frobenius sends (\Theta) to (\Theta\^{}2=1), and (\Theta-1\ne0) in
(\mathbb F\_2{[}\Theta{]}/(\Theta-1)\^{}2). This proves nonexistence.
Consequently the full product
(B\_\{\mathbb Z\_2\}\times I\_\{\mathbb Z\_2\}) admits no Frobenius
lift, by the idempotent argument just proved, and hence no delta
structure on this torsionfree ring. At odd primes the infinity formula
(\Phi(\Theta)=\Theta\^{}p) respects (Theta\^{}2=-1) and lifts Frobenius,
so it supplies a delta structure and a bounded crystalline prism. In
particular the full rank-eight model at3, using the augmentation lift on
(B\_K), is a valid bounded crystalline prism and has both classes of
(IEC43).

\subsection{9. Support and the scope of the
comparison}\label{support-and-the-scope-of-the-comparison}

Every displayed integral ring map induces the programme's
support-preserving semiring map (G\_L(f)), sending a supported amplitude
((r,1)) to ((f(r),1)) and each zero-fibre point ((0,\lambda)) to itself.
Multiplication and addition are preserved because (f) is a ring map and
the support meet/join coordinates do not change. In particular the
quotient (C\_K\to D\_K) sends the nonzero amplitude (T) to a nonzero
amplitude (eta) whose square is the supported zero (e\_L=(0,1)), not the
external element ( au\_L=(0,0)). Its kernel consists of the actual
amplitudes (K\epsilon\oplus K\epsilon T); these map to supported zero.
The prime-primary classes and their maps above therefore have defined
supported receivers without identifying ( au\_L) with a ring zero or
deleting (e\_L).

The trace maps (IEC12)--(IEC15), the exact cotangent maps
(IEC19)--(IEC22), (IEC45)--(IEC47), and the cohomology maps (IEC40),
(IEC44), (IEC50), (IEC55) specify different observations of these
receivers. The PD comparison retains a nonzero class at2 and at3, while
the integral lattice adds further classes with their stated kernels. Any
application to a supported-zero-prime Weil functional must use that
functional's actual formula and its source map; these integral
calculations make no substitution of an older trace for that missing
analytic identification.

\subsection{Appendix A. Complete independent proof for the unsplit
prime3
component}\label{appendix-a.-complete-independent-proof-for-the-unsplit-prime3-component}

The following proof is retained in full, including its original P3
labels, so all claims imported in Section7 have their derivation in this
document. \# The full integral collision at the prime 3

This calculation retains the finite collision algebra, its original root
coordinate, and the signed coordinate. All coefficient maps fix
\(O=\mathbb Z_3\). Put \[
B=O[R,T]/((R-1)^2(R+2),\ T^2-3R^2+3),\qquad x=R-1.
\tag{P3-1}
\] The substitution is an isomorphism of the displayed presentations,
with inverse \(R=x+1\). Thus \[
B=O[x,T]/(f,g),\quad f=x^2(x+3),\quad
g=T^2-3x(x+2),\quad D=O[T]/T^2,
\quad\pi:B\longrightarrow D,\quad(x,T)\longmapsto(0,T).
\tag{P3-2}
\] The reflection is \(\sigma(x)=x,\ \sigma(T)=-T\), on both rings.
Every calculation below takes place over \(O\); in particular, no root
projector involving \(1/3\) is used to define an integral map.

The incoming coordinate identities are
\href{https://github.com/KokunoYumeto/zeta-function-research-reader/blob/d5c9d198a8432e3ede468b280162a99c00e3f4f7/workbenches/splitzero-tandem/branches/tau-zero-prime-positivity/EIGHT_STATE_HEAT_COMPARISON.md}{EIGHT\_STATE\_HEAT\_COMPARISON.md,
ESH58--72}. Those statements were read directly together with the
relevant entries in SOURCE\_READING\_AND\_USE.md. The Frobenius/delta
definition used here is Bhargav Bhatt and Peter Scholze, \emph{Prisms
and Prismatic Cohomology},
\href{https://arxiv.org/src/1905.08229v4}{arXiv:1905.08229v4, original
author source}, \texttt{prisms.tex} lines 484--541, in particular
\texttt{def:deltaring} and \texttt{FrobLiftTheta}, read directly for
this calculation. The cotangent foundation is Luc Illusie,
\emph{Complexe cotangent et déformations I}, III, 3.2.2--3.2.4, in the
\href{https://zenodo.org/records/22884471}{retained edition}; the
diplomatic TeX lines 23386--23779, including the proof of 3.2.4, were
read directly. That edition is a transcription of the historical work,
not original author TeX. Its source identity and the original
Bhatt--Scholze archive hashes remain those recorded in
SOURCE\_READING\_AND\_USE.md. All specialized maps, lattices, kernel
computations, and cohomology groups below are derived explicitly here.

\subsection{1. The retained algebra is local and
free}\label{the-retained-algebra-is-local-and-free}

\textbf{P3.1.} The ordered list \[
\mathcal E=(1,x,x^2,T,xT,x^2T)
\tag{P3-3}
\] is an \(O\)-basis of \(B\). The ring \(B\) is complete, has no
3-torsion, and is local with maximal ideal \((3,x,T)\). In particular,
it admits no nontrivial decomposition as a product of rings over \(O\).

\textbf{Proof.} Division by the monic polynomial \(f\) gives the free
rank-three ring \(K=O[x]/f\). Division by the monic polynomial
\(T^2-3x(x+2)\) then gives \(B=K\oplus TK\), proving the asserted basis
and absence of 3-torsion. A finite free module over the complete ring
\(O\) is complete. Reduction modulo 3 is \[
B/3B=\mathbb F_3[x,T]/(x^3,T^2),
\tag{P3-4}
\] which has the unique maximal ideal \((x,T)\). Because \(B\) is finite
and integral over \(O\), every maximal ideal of \(B\) contracts to
\((3)\); consequently (P3-4) identifies its unique maximal ideal. An
idempotent in a local ring is 0 or 1: one of \(e\) and \(1-e\) is a
unit, and \(e(1-e)=0\) then forces the other to vanish. A nontrivial
product would supply a different idempotent. This proves the last
assertion. \(\square\)

\subsection{2. The integral nilradical, its powers, and the reduced
order}\label{the-integral-nilradical-its-powers-and-the-reduced-order}

Retain exactly the three generators from the incoming collision
calculation: \[
n_1=x(x+3),\qquad n_2=T(x+3),\qquad n_3=Tx(x+3).
\tag{P3-5}
\]

\textbf{P3.2.} The nilradical and all its nonzero powers are \[
\begin{aligned}
N&=On_1\oplus On_2\oplus On_3,\\
N^2&=18On_1\oplus3On_3,\\
N^3&=54On_3,\qquad N^4=0.
\end{aligned}
\tag{P3-6}
\] The complete multiplication and the actions of the algebra generators
are \[
\begin{gathered}
n_1^2=n_1n_3=n_2n_3=n_3^2=0,\qquad
n_2^2=18n_1,\qquad n_1n_2=3n_3,\\
xn_1=xn_3=0,\quad xn_2=n_3,\qquad
Tn_1=n_3,\quad Tn_2=6n_1,\quad Tn_3=0.
\end{gathered}
\tag{P3-7}
\]

\textbf{Proof.} The defining equations give \(x^3=-3x^2\), \(x^4=9x^2\),
and \(T^2=3x(x+2)\). Hence \[
n_1^2=x^2(x+3)^2=f(x+3)=0,\qquad
xn_1=f=0,
\] and \[
Tn_2=3x(x+2)(x+3)=3(x+2)n_1=6n_1.
\] Multiplication of these identities by \(T\), together with
\(xn_2=n_3\), proves the second row of (P3-7). The remaining products
are \[
n_2^2=(x+3)Tn_2=6(x+3)n_1=18n_1,
\quad n_1n_2=(x+3)Tn_1=3n_3,
\] and the zero products follow from \(xn_1=0\), \(n_1^2=0\), and
\(Tn_3=0\). Thus their span is an ideal of fourth power zero. Its three
displayed generators are independent in (P3-3).

There are three \(O\)-algebra maps to \(O\), with values on \((x,T)\)
equal to \[
(0,0),\qquad(-3,3),\qquad(-3,-3).
\tag{P3-8}
\] For \(p,q\in O[x]\) of degree at most two, their simultaneous
vanishing on \(p+Tq\) means \[
p(0)=0,\qquad p(-3)+3q(-3)=0,\qquad
p(-3)-3q(-3)=0.
\] Since 2 and 3 are nonzero divisors in \(O\), this is equivalent to
\(p(0)=p(-3)=q(-3)=0\). Writing coefficients explicitly gives
\(p=a(x^2+3x)\) and \(q=b(x+3)+c(x^2+3x)\), with \(a,b,c\in O\).
Therefore the common kernel of (P3-8) is exactly the stated span. Every
nilpotent must vanish under (P3-8), because \(O\) has no nonzero
nilpotents. The span already has fourth power zero, so it is the entire
nilradical.

The nonzero entries of the multiplication table generate \(N^2\).
Multiplying these by \(N\) gives only multiples of \(54n_3\), with
\((18n_1)n_2=54n_3\); multiplying once more gives zero. Independence of
\(n_1,n_3\) proves that the sums and powers in (P3-6) are exact.
\(\square\)

\textbf{P3.3.} Evaluation in (P3-8) induces \[
B/N\ \simeq\ \{(a,b,c)\in O^3:a\equiv b\equiv c\pmod3\}.
\tag{P3-9}
\] The integral closure of this reduced order in \(\mathbb Q_3^3\) is
\(O^3\); its conductor is \(3O^3\), and \[
O^3/(B/N)\simeq (O/3O)^2.
\tag{P3-10}
\]

\textbf{Proof.} Modulo \(N\), the basis can be reduced to \(1,x,T\),
using \[
x^2=-3x,\quad xT=-3T,\quad x^2T=9T,\quad T^2=-3x.
\] The images of these three basis vectors under (P3-8) are \((1,1,1)\),
\((0,-3,-3)\), and \((0,3,-3)\). Their span lies in the right side of
(P3-9). Conversely a triple there has coefficients \[
a\cdot1+\frac{2a-b-c}{6}x+\frac{b-c}{6}T.
\] Both quotients belong to \(O\), since their numerators lie in \(3O\)
and 2 is a unit; the formula asserts divisibility within \(O\), not
inversion of 3 in the algebra. This proves (P3-9). The two independent
differences modulo 3 give (P3-10).

The ring \(O^3\) is integral over \(B/N\), because it is generated by
the coordinate idempotents and each satisfies \(z^2-z=0\). It is
integrally closed in \(\mathbb Q_3^3\): a rational 3-adic coordinate
integral over \(O\) has nonnegative valuation. Conversely an element of
\(\mathbb Q_3^3\) integral over \(B/N\) is integral over \(O\) by
transitivity, because \(B/N\) is finite over \(O\). Thus every
coordinate belongs to \(O\), proving the integral-closure assertion. An
element \(z\in O^3\) lies in the conductor exactly when each \(ze_i\)
lies in (P3-9), where \(e_i\) are the coordinate idempotents. This is
equivalent to all three coordinates of \(z\) belonging to \(3O\).
\(\square\)

\subsection{3. Exact defects of the marked-root
map}\label{exact-defects-of-the-marked-root-map}

\textbf{P3.4.} The map \(\pi:B\to D\) is onto, with kernel \(xB\). Its
restriction to the nilradical satisfies \[
\pi(n_1)=\pi(n_3)=0,\qquad\pi(n_2)=3T,
\tag{P3-11}
\] and consequently there is an exact sequence of \(O\)-modules \[
0\longrightarrow On_1\oplus On_3\longrightarrow N
\xrightarrow{\ \pi\ }OT\longrightarrow (O/3O)T\longrightarrow0.
\tag{P3-12}
\] The comparison with the square of the nilradical is \[
\begin{aligned}
\ker(\pi|_N)/N^2&\simeq (O/9O)n_1\oplus(O/3O)n_3,\\
N/N^2&\simeq On_2\oplus(O/9O)n_1\oplus(O/3O)n_3,\\
N^2/N^3&\simeq O(18n_1)\oplus(O/9O)(3n_3).
\end{aligned}
\tag{P3-13}
\] All these maps preserve reflection; \(n_1\) has sign \(+1\), and
\(n_2,n_3,T\) have sign \(-1\).

\textbf{Proof.} Setting \(x=0\) in the presentation leaves exactly
\(O[T]/T^2\), proving the first assertion. Substitution proves (P3-11).
Independence of the generators proves the kernel and image in (P3-12),
and (P3-6) gives every quotient in (P3-13), using that 18 generates the
ideal \(9O\) and 54 generates \(27O\). The reflection signs follow
directly from (P3-5). \(\square\)

\textbf{P3.5.} The square-zero elements of \(B\) are exactly
\(On_1\oplus On_3\). Every \(O\)-algebra map \(j:D\to B\) therefore
satisfies \(\pi j(T)=0\). In particular \(\pi\) has no algebra section.

\textbf{Proof.} A square-zero element is nilpotent, so it has a unique
expression \(v=an_1+bn_2+cn_3\). The table gives \[
v^2=18b^2n_1+6abn_3.
\tag{P3-14}
\] This is zero exactly when \(b=0\), because \(O\) is a domain and the
generators are independent. The image of \(T\) under \(j\) must be
square-zero; (P3-11) proves its zero marked-root image. \(\square\)

For an additional exact comparison, retain the integral double-root
algebra \[
C=O[\epsilon,T]/(\epsilon^2,T^2-6\epsilon).
\tag{P3-15}
\]

\textbf{P3.6.} The coordinate-preserving map \[
\iota:B\longrightarrow C\times O\times O,
\quad x\longmapsto(\epsilon,-3,-3),\quad
T\longmapsto(T,3,-3)
\tag{P3-16}
\] is injective, and its cokernel as an \(O\)-module is \[
\operatorname{coker}\iota\simeq O/9O\oplus O/27O.
\tag{P3-17}
\]

\textbf{Proof.} Both defining equations vanish after the indicated
substitutions. In the ordered basis \(1,\epsilon,T,\epsilon T\) of
\(C\), an element of the target is written
\((a+b\epsilon+cT+d\epsilon T,e,f)\). It lies in the image if and only
if \[
\frac{e+f}{2}-a+3b\in9O,\qquad
\frac{e-f}{2}-3c+9d\in27O.
\tag{P3-18}
\] Indeed a source element \(p_0+p_1x+p_2x^2+T(q_0+q_1x+q_2x^2)\) has
\(a=p_0,b=p_1,c=q_0,d=q_1\); the two expressions in (P3-18) equal
\(9p_2\) and \(27q_2\). Conversely membership in the respective ideals
determines unique \(p_2,q_2\in O\), and these produce the proposed
target. This also proves injectivity. Reduction of the two expressions
modulo their indicated ideals is onto \(O/9O\oplus O/27O\), since the
coordinates \(e,f\) allow arbitrary sum and difference after division by
the unit 2. Its kernel is exactly the image just computed, proving
(P3-17). \(\square\)

\subsection{4. Every coefficient-preserving Frobenius
lift}\label{every-coefficient-preserving-frobenius-lift}

Here a Frobenius lift means an \(O\)-algebra endomorphism whose
reduction modulo 3 is the absolute Frobenius of (P3-4).

\textbf{P3.7.} Every such endomorphism is exactly one of the following:
\[
\begin{array}{ll}
\Phi_+:&(x,T)\longmapsto(-3,3),\\
\Phi_-:&(x,T)\longmapsto(-3,-3),\\
\Phi_v:&x\longmapsto\dfrac32v^2,\quad T\longmapsto3v,
\qquad v\in N.
\end{array}
\tag{P3-19}
\] The parameter \(v\) in the third family is unique. For
\(v=an_1+bn_2+cn_3\), its coordinates are \[
\Phi_v(x)=27b^2n_1+9abn_3,\qquad
\Phi_v(T)=3an_1+3bn_2+3cn_3.
\tag{P3-20}
\]

\textbf{Proof.} In (P3-4), \(x^3=T^3=0\), so every Frobenius lift has
images \(X=3u,Y=3v\), for unique \(u,v\in B\). Substituting these into
the defining equations and cancelling powers of the nonzero divisor 3
gives the equivalent conditions \[
u^2(u+1)=0,\qquad v^2=u(3u+2).
\tag{P3-21}
\] Because \(B\) is local, either \(u\) is a unit or \(u+1\) is a unit.
In the first case (P3-21) implies \(u=-1\), then \(v^2=1\). Since
\((v-1)(v+1)=0\) and their difference is the unit 2, one factor is a
unit and the other is zero. Thus \(v=1\) or \(v=-1\), producing the
first two maps.

If \(u\) is not a unit, then \(u+1\) is a unit, so \(u^2=0\). The other
equation gives \(v^2=2u\), hence \(v^4=0\) and \(v\in N\). It follows
that \(u=v^2/2\), yielding the third family. Conversely each \(v\in N\)
has fourth power zero by (P3-6); therefore \(u=v^2/2\) satisfies both
equations in (P3-21). All listed images are divisible by 3, so they lift
Frobenius. Uniqueness follows from \(\Phi_v(T)=3v\) and the absence of
3-torsion. Formula (P3-20) is (P3-14) substituted into (P3-19).
\(\square\)

\textbf{P3.8.} The delta structures belonging to these lifts have the
exact generator values \[
\begin{aligned}
\delta_v(x)&=x^2+\frac12v^2,&
\delta_v(T)&=v-Tx(x+2),\\
\delta_\pm(x)&=x^2-1,&
\delta_\pm(T)&=\pm1-Tx(x+2),\\
\delta(c)&=(c-c^3)/3&& (c\in O).
\end{aligned}
\tag{P3-22}
\] In particular the augmentation \(\Phi_0(x)=\Phi_0(T)=0\) is a
reflection-equivariant Frobenius lift on the whole unsplit algebra, with
\[
\delta_0(x)=x^2,\qquad \delta_0(T)=-Tx(x+2).
\tag{P3-23}
\]

\textbf{Proof.} The differences \(\Phi(z)-z^3\) lie in \(3B\) by the
Frobenius condition, and division by 3 is unique because \(B\) has no
3-torsion. Expanding \(\Phi(a+b)=\Phi(a)+\Phi(b)\) and
\(\Phi(ab)=\Phi(a)\Phi(b)\) gives directly \[
\delta(a+b)=\delta(a)+\delta(b)-a^2b-ab^2,
\quad
\delta(ab)=a^3\delta(b)+b^3\delta(a)+3\delta(a)\delta(b),
\] with \(\delta(0)=\delta(1)=0\). Thus these are delta structures. Now
\(x^3=-3x^2\) and \(T^3=3Tx(x+2)\), proving (P3-22). Augmentation and
reflection commute on both generators, and (P3-23) is the case \(v=0\).
\(\square\)

\textbf{P3.9.} The Frobenius lifts commuting with reflection are exactly
\[
\Phi_{bn_2+cn_3}\qquad(b,c\in O).
\tag{P3-24}
\] The maps in (P3-19) that induce a Frobenius lift on \(D\) through
\(\pi\) are exactly the family \(\Phi_v\); their induced maps are \[
\Phi_D(T)=9bT,\qquad\delta_D(T)=3bT
\quad\text{for }v=an_1+bn_2+cn_3.
\tag{P3-25}
\] Every coefficient-preserving Frobenius lift on \(D\) has the form
\(T\mapsto3cT\), \(c\in O\), and it comes from a lift on \(B\) if and
only if \(c\in3O\).

\textbf{Proof.} For \(\Phi_v\), commuting on \(T\) is equivalent to
\(\sigma(v)=-v\). In the independent basis (P3-5), this is precisely
\(2a=0\), hence \(a=0\). When it holds, \(v^2\) is fixed by reflection,
so commuting on \(x\) also holds. The constant maps \(\Phi_\pm\) cannot
commute on \(T\), since they send it to a nonzero fixed constant. This
proves (P3-24).

The kernel of \(\pi\) is \(xB\). For \(\Phi_v\), equations (P3-11) and
(P3-14) show that \(\pi\Phi_v(x)=0\); the induced value of \(T\) is
\(3\pi(v)=9bT\). In contrast \(\pi\Phi_\pm(x)=-3\ne0\), so the kernel is
not preserved. Since both rings have no 3-torsion, a commuting Frobenius
square also commutes with the associated delta maps, proving (P3-25).

If an endomorphism of \(D\) sends \(T\) to \(a+bT\), its square-zero
condition gives \(a^2=0\), so \(a=0\). Its reduction is Frobenius
exactly when \(b\in3O\). By (P3-25), liftability through \(\pi\)
requires \(b\in9O\), and every value \(b=9d\) is attained by choosing
\(v=dn_2\). \(\square\)

\subsection{5. The entire cotangent
presentation}\label{the-entire-cotangent-presentation}

\textbf{P3.10.} The relative cotangent complex is the following two-term
complex, in cohomological degrees \(-1,0\): \[
L_{B/O}\simeq
\left[B e_f\oplus B e_g
\xrightarrow{J}B\,dx\oplus B\,dT\right],
\qquad
J=\begin{pmatrix}
3x(x+2)&-6(x+1)\\0&2T
\end{pmatrix}
=\begin{pmatrix}T^2&-6R\\0&2T\end{pmatrix}.
\tag{P3-26}
\] It has no cohomology outside those degrees. Its degree \(-1\)
cohomology is free over \(O\) on \[
u_1=(n_1,0),\qquad u_2=(n_3,0),\qquad
u_3=(n_2,n_3).
\tag{P3-27}
\] There is an explicit isomorphism of \(B\)-modules \[
N\xrightarrow{\sim}H^{-1}(L_{B/O}),
\qquad n_1\mapsto u_1,\quad n_3\mapsto u_2,\quad n_2\mapsto u_3.
\tag{P3-28}
\]

\textbf{Proof.} In \(P=O[x,T]\), \(f\) is a nonzero divisor. The ring
\(P/(f)\) is \(K[T]\), and \(g\) is monic in \(T\), so multiplication by
\(g\) is injective: the leading coefficient of a nonzero polynomial
remains the leading coefficient of its product with \(g\). Thus
\((f,g)\) is a regular sequence. For this regular immersion, the
conormal module is free on \(e_f,e_g\) and the regular-immersion
cotangent description, together with the polynomial base, gives the
two-term complex. Differentiation gives exactly (P3-26).

For completeness, the needed annihilators can be solved in the free
basis. In \(K\), \[
x(a+bx+cx^2)=ax+(b-3c)x^2,
\] so \(\operatorname{ann}_K(x)=On_1\). The element \(x+2\) is a unit
because its residue is 2 in the residue field; therefore
\(\operatorname{ann}_K(3x(x+2))=On_1\). Writing an element of \(B\) as
\(p+Tq\), the equation \(T(p+Tq)=0\) first forces \(p=0\) and then
\(q\in On_1\). Consequently \[
\operatorname{ann}_B(T)=On_3,\qquad
\operatorname{ann}_B(T^2)=On_1\oplus On_3.
\tag{P3-29}
\] If \(J(a,b)=0\), its second coordinate gives \(b=cn_3\) for a unique
\(c\in O\). Since \(Rn_3=n_3\) and \(T^2n_2=6n_3\), the first coordinate
becomes \(T^2(a-cn_2)=0\). Equation (P3-29) now gives the unique
expression \[
(a,b)=Au_1+Cu_2+cu_3\qquad(A,C,c\in O).
\] This proves (P3-27). The actions on this basis are \[
xu_1=xu_2=0,\quad xu_3=u_2,\qquad
Tu_1=u_2,\quad Tu_2=0,\quad Tu_3=6u_1.
\] They agree term by term with (P3-7) under the correspondence in
(P3-28), proving that it is \(B\)-linear. \(\square\)

\textbf{P3.11.} The degree-zero cotangent cohomology has the following
exact decomposition as an \(O\)-module: \[
\Omega^1_{B/O}
=Oa_0\oplus Ob_0\oplus Ob_1
\oplus(O/3O)a_1\oplus(O/3O)a_2
\oplus(O/3O)a_4\oplus(O/3O)a_5
\oplus(O/3O)b_2\oplus(O/3O)c,
\tag{P3-30}
\] where every symbol denotes the following actual form: \[
\begin{gathered}
a_0=dx,\quad a_1=x\,dx,\quad a_2=x^2\,dx,\quad
a_3=T\,dx,\quad a_4=xT\,dx,\quad a_5=x^2T\,dx,\\
b_0=dT,\quad b_1=x\,dT,\quad b_2=x^2\,dT,\quad
c=a_3-2b_1=T\,dx-2x\,dT.
\end{gathered}
\tag{P3-31}
\] In particular \(\Omega^1_{B/O}\simeq O^3\oplus(O/3O)^6\).

\textbf{Proof.} Before imposing relations, also put \(b_3=T\,dT\),
\(b_4=xT\,dT\), \(b_5=x^2T\,dT\). Multiply \(df\) by the six basis
elements (P3-3). The resulting relations are equivalent to \[
3a_1=3a_2=3a_4=3a_5=0:
\] the even ones are \(6a_1+3a_2=0\), \(-3a_2=0\), \(9a_2=0\); the odd
ones are \(6a_4+3a_5=0\), \(-3a_5=0\), \(9a_5=0\). These equivalences
use only division by the unit 2.

Multiplication of \(dg=-6(x+1)dx+2T\,dT\) by \(1,x,x^2\) then
eliminates, with coefficient 1, the three symbols \[
b_3=3(a_0+a_1),\qquad
b_4=3(a_1+a_2)=0,\qquad b_5=-6a_2=0.
\] Its products by \(T,xT,x^2T\) give respectively \[
3(-a_3-a_4+2b_1+b_2)=0,\quad
3(a_4+a_5+b_2)=0,\quad
6a_5+9b_2=0.
\] With the already derived relations these are equivalent to \(3b_2=0\)
and \(3(a_3-2b_1)=0\). Every generator of the original relation module
has been included, and every elimination used an invertible coefficient.
Replacing \(a_3\) by \(c+2b_1\) therefore gives precisely (P3-30), with
no additional relations. \(\square\)

\subsection{6. Ordinary de Rham cohomology, with
representatives}\label{ordinary-de-rham-cohomology-with-representatives}

Here the complex is the ordinary relative Kähler de Rham complex
\(\Omega^\bullet_{B/O}\). Its calculation does not identify this complex
with crystalline or prismatic cohomology.

\textbf{P3.12.} Its terms in degree at least two are \[
\Omega^2_{B/O}=\mathbb F_3[x]/(x^3)\,w,
\qquad w=dx\wedge dT,\qquad \Omega^j_{B/O}=0\quad(j\ge3).
\tag{P3-32}
\] The entire ordinary de Rham cohomology is \[
\begin{aligned}
H^0_{\mathrm{dR}}(B/O)&=O\cdot1\oplus O\cdot(3x^2)
\oplus O\cdot(3x^2T),\\
H^1_{\mathrm{dR}}(B/O)&=(O/3O)[x^2dx]
\oplus(O/3O)[3x\,dT],\\
H^j_{\mathrm{dR}}(B/O)&=0\qquad(j\ge2).
\end{aligned}
\tag{P3-33}
\] Both displayed degree-one classes are nonzero and have exact order 3.

\textbf{Proof.} Wedge the two relations in (P3-26) with \(dx,dT\). Their
coefficient ideal on \(w\) is \((T^2,6R,2T)\). The elements 2 and
\(R=x+1\) are units, so this ideal is \((3,T)\). Taking the quotient in
(P3-4) gives (P3-32). Exterior powers of a module generated by two
elements vanish in degrees at least three.

Using the direct decomposition (P3-30), the differential from the six
basis vectors of \(B\) is \[
\begin{array}{c|rrrrrr}
z&1&x&x^2&T&xT&x^2T\\ \hline
dz&0&a_0&2a_1&b_0&c+3b_1&2a_4+b_2.
\end{array}
\tag{P3-34}
\] For \(z=p_0+p_1x+p_2x^2+q_0T+q_1xT+q_2x^2T\), the free coordinates of
\(dz\) force \(p_1=q_0=q_1=0\). Its remaining torsion coordinates vanish
exactly when \(p_2,q_2\in3O\), proving the first line of (P3-33).

The next differential, on the generators of (P3-30), is \[
\begin{gathered}
da_0=db_0=da_1=da_2=dc=0,\qquad db_1=w,\\
da_4=-xw,\quad da_5=-x^2w,\quad db_2=2xw.
\end{gathered}
\tag{P3-35}
\] For instance \(dc=d(Tdx)-2d(xdT)=-3w=0\). These formulas reach all
three basis vectors \(w,xw,x^2w\) of (P3-32), proving \(H^2=0\).

Put \(u=2a_4+b_2\). Since \(w,xw,x^2w\) are independent over
\(\mathbb F_3\), (P3-35) gives the full degree-one kernel as \[
Oa_0\oplus Ob_0\oplus O(3b_1)
\oplus(O/3O)a_1\oplus(O/3O)a_2
\oplus(O/3O)c\oplus(O/3O)u.
\tag{P3-36}
\] Indeed the coefficient of \(b_1\) must be divisible by 3; the
coefficient of \(a_5\) must vanish; and the relation between the
coefficients of \(a_4,b_2\) is exactly their being a multiple of
\((2,1)\). Now (P3-34) says that the image from degree zero is generated
by \[
a_0,\quad b_0,\quad 2a_1,\quad c+3b_1,\quad u.
\] After quotienting, \(a_2\) remains an independent order-three
generator. If \(s=3b_1\), the other remaining summand is \[
(Os\oplus(O/3O)c)/(s+c)\simeq(O/3O)s.
\] This proves the second line of (P3-33), including nonvanishing and
exact orders. Higher vanishing follows from (P3-32). \(\square\)

The closed degree-zero elements form an explicit order. With \(A=3x^2\)
and \(E=3x^2T\), their multiplication is \[
A^2=27A,\qquad AE=27E,\qquad E^2=243A.
\tag{P3-37}
\] Indeed \(x^4=9x^2\) gives the first two identities, while
\(x^4T^2=81x^2\) gives the last. Under the three evaluations (P3-8),
\(A\) has values \((0,27,27)\) and \(E\) has values \((0,81,-81)\). Thus
the degree-zero calculation keeps the original integral constants and
their exact indices.

\subsection{7. The cotangent and de Rham maps to the dual
numbers}\label{the-cotangent-and-de-rham-maps-to-the-dual-numbers}

\textbf{P3.13.} The cotangent map induced by \(\pi\) is represented, in
the displayed polynomial presentations, by \[
\begin{array}{rcl}
B e_f\oplus B e_g&\longrightarrow&D e_{T^2},
\qquad(a,b)\longmapsto\pi(b)e_{T^2},\\
B\,dx\oplus B\,dT&\longrightarrow&D\,dT,
\qquad a\,dx+b\,dT\longmapsto\pi(b)\,dT,
\end{array}
\tag{P3-38}
\] where \(L_{D/O}=[D e_{T^2}\xrightarrow{2T}D\,dT]\). It induces the
zero map on \(H^{-1}\). On degree zero it sends \(b_0\mapsto dT\) and
every other generator in (P3-30) to zero. The ordinary de Rham map is
augmentation on \(H^0\), and both classes of \(H^1_{\mathrm{dR}}(B/O)\)
have zero image.

\textbf{Proof.} The map of polynomial rings sends \(f\mapsto0\) and
\(g\mapsto T^2\). Therefore its map on conormal generators is
\(e_f\mapsto0,e_g\mapsto e_{T^2}\); the map on differentials is
\(dx\mapsto0,dT\mapsto dT\). This proves (P3-38) and also verifies the
chain-map identity directly with (P3-26). Each cycle in (P3-27) has
second coordinate either 0 or \(n_3\), whose image is zero by (P3-11),
proving the \(H^{-1}\) assertion. Substitution gives the claimed images
in degree zero.

For \(D\), the relation is \(2T\,dT=0\). Since 2 is a unit,
\(\Omega^1_{D/O}=O\,dT\), and the derivative \(a+bT\mapsto b\,dT\) is
onto with kernel \(O\). Hence \(H^0_{\mathrm{dR}}(D/O)=O\) and
\(H^j_{\mathrm{dR}}(D/O)=0\) for \(j>0\). Applying \(x\mapsto0\) to
(P3-33) proves all de Rham assertions. \(\square\)

\textbf{P3.14.} Every Frobenius lift in (P3-19) acts as zero on
\(H^1_{\mathrm{dR}}(B/O)\). For \(\Phi_v\), its action on \(H^0\) is
augmentation; the actions of \(\Phi_\pm\) on the basis \(1,A,E\) of
(P3-37) are \[
(1,A,E)\longmapsto(1,27,\pm81).
\tag{P3-39}
\]

\textbf{Proof.} For \(\Phi_v\), put \(X=3v^2/2,Y=3v\). Then \(X^2=0\),
so \(\Phi_v(A)=\Phi_v(E)=0\). On the two degree-one representatives its
pullback gives \[
x^2dx\longmapsto X^2dX=0,\qquad
3x\,dT\longmapsto3X\,dY
=\frac{27}{2}v^2dv=\frac92d(v^3).
\] The last expression is exact over \(O\), since \(9/2\in O\). The
constant maps kill every positive-degree differential and have the
values (P3-39) by direct substitution. \(\square\)

The results above distinguish three explicit maps on the retained
infinitesimal: \(N\to OT\) has image \(3OT\); the induced map on degree
\(-1\) cotangent cohomology is zero; and the two nonzero ordinary de
Rham classes in (P3-33) map to zero. Equations (P3-11), (P3-27), and
(P3-38) prove each connecting map in the original coordinates.

\end{document}
